% ===================================================================
%  TABELL Nr 16 — Det stora varuhuset. — Basaren. — Leksakerna.
%  Traduction 2026 d'apres texto/io, controlee sur texto/fr et
%  texto/nl. Consigne : voir texto/sv/10-tabell-01.tex.
%
%  Le tableau d'astronomie (bloc 15-1) porte des renvois a LETTRES,
%  (a) a (f), graves sur le tableau lui-meme. L'ido saute la lettre
%  (d), que le fac-simile grave pourtant : la colonne la rend, comme
%  toutes les autres traductions. Ecart declare dans outils/renvoji.py.
%
%  DERNIER TABLEAU DE LA COLONNE SUEDOISE. Bilan des seize : une
%  seule inversion vraie, au tableau 7, et trois evitees d'avance aux
%  tableaux 5, 14 et 15. Toutes tenaient au meme trait — le genitif
%  prepose — et toutes se defont par la meme voie : une preposition a
%  la place du 's', ou un complement porte en tete de phrase, que la
%  seconde position du verbe rend gratuite.
% ===================================================================

%%K t16-apar-1 apar
\VUsaut{4.0mm}
\VUtitre{15.0pt}{60}{TABELL Nr 16}
\VUsaut{2.0mm}
\VUpk{13.2pt}{40}{Det stora varuhuset. --- Basaren.}
\VUpk{13.2pt}{40}{Leksakerna.}
\VUsaut{3.4mm}

%%K t16-01-1 p
1. --- Det nya året nalkas. De stora varuhusen har gjort i ordning sina utställningar av
\VUgras{leksaker} \textsuperscript{(1)} och nyårsgåvor.

%%K t16-01-2 p
Jag frågade min farfar, om han ville föra mig till basaren för att se den vackra
utställningen, och han samtyckte till det.

%%K t16-02-1 p
2. --- Först stannade vi framför vackra \VUgras{vapensköldar}, som innehöll ett
\VUgras{gevär}, en \VUgras{käppi}, en \VUgras{sabel}, ett \VUgras{värjbälte} och
ett par \VUgras{epåletter}, eller en \VUgras{hjälm}, ett \VUgras{kyrass} \textsuperscript{(3)}
och en \VUgras{pistol} \textsuperscript{(4)}. Allt det intresserade mig mycket mer än
\VUgras{trollyktorna} \textsuperscript{(2)}.

%%K t16-tit-1 sub
\VUsaut{3.0mm}
\VUpk{10.2pt}{40}{Vackra syner.}
\VUsaut{2.6mm}

%%K t16-03-1 p
3. --- I samma avdelning stod en \VUgras{skiltvakt} \textsuperscript{(5)} i sin
\VUgras{vaktkur}; han tycktes hålla vakt framför ett helt menageri av papp. Hur
många djur fanns där inte: en \VUgras{papegoja} \textsuperscript{(6)} på sin pinne, en
\VUgras{noshörning} \textsuperscript{(7)} med ett horn på näsan, en
\VUgras{ren} \textsuperscript{(8)}, en \VUgras{struts} \textsuperscript{(9)}, en
\VUgras{apa} \textsuperscript{(10)} som knäckte en nöt, en \VUgras{räv}
\textsuperscript{(11)} som bar bort en höna, en \VUgras{krokodil} \textsuperscript{(12)} som
gapade med väldiga käftar, en \VUgras{kamel} \textsuperscript{(13)}, en graciös
\VUgras{vinthund} \textsuperscript{(14)}, en
\VUgras{panter} \textsuperscript{(15)}, sedan två djur som jag inte kände till
och som, säger man, är en \VUgras{vessla} \textsuperscript{(16)} och en
\VUgras{hyena} \textsuperscript{(17)}. Där fanns också en \VUgras{leopard}
\textsuperscript{(18)} och en \VUgras{varg} \textsuperscript{(21)}; alldeles intill fanns
näpna \VUgras{råttor} \textsuperscript{(19)} och pyttesmå
\VUgras{möss} \textsuperscript{(20)} av gummi.

%%K t16-03-1 p suite
Till slut avslutade en \VUgras{flodhäst} \textsuperscript{(22)} och en
\VUgras{dromedar} \textsuperscript{(23)} med en puckel samlingen. Jag skulle ha
stannat där flera timmar till om det inte alldeles intill hade funnits ett praktfullt läger
fullt av \VUgras{tennsoldater} \textsuperscript{(29)} som drog till sig min uppmärksamhet.

%%K t16-tit-2 sub
\VUsaut{4.0mm}
\VUpk{10.2pt}{40}{Soldater och krig.}
\VUsaut{2.8mm}

%%K t16-04-1 p
4. --- Min farfar, som är \VUgras{invalid} \textsuperscript{(24)} och som måste lämna armén på
grund av ett \VUgras{sår} varigenom han fick
\VUgras{krigsmedaljen}, berättar gärna för mig om de \VUgras{fälttåg} som han
deltog i. Han var lycklig medan han förklarade den lilla arméns olika rörelser i miniatyr. En
grupp riktiga soldater som besökte basaren: en \VUgras{ingenjörsoldat}
\textsuperscript{(25)}, en \VUgras{alpjägare} \textsuperscript{(26)} med huvudet täckt av en
\VUgras{basker}, en \VUgras{fanjunkare} \textsuperscript{(27)} vid infanteriet och en
\VUgras{sergeant vid artilleriet} \textsuperscript{(28)} med en \VUgras{trens} av guld på
ärmen, stannade bredvid oss för att lyssna till förklaringarna.

%%K t16-05-1 p
5. --- Min farfar, mycket stolt över att ha en så lysande åhörarskara, improviserade en
fullständig slagplan.

%%K t16-05-2 p
Den befästa staden, med sina \VUgras{vallar} och \VUgras{gravar}, belägrades av den
\VUgras{fientliga armén}, som, efter ett långt \VUgras{bombardemang}, var redo till
\VUgras{stormning}. Men de \VUgras{belägrade} som var tvingade av hungern, hade vågat ett
\VUgras{utfall} mot \VUgras{lägret} hos de
\VUgras{belägrande}. Sedan de hade slagit tillbaka \VUgras{anfallet} med en
envis \VUgras{strid} runt \VUgras{tältlägret}, sände de \VUgras{segrande} belägrarna ut sina
\VUgras{skvadroner} för att förfölja de \VUgras{besegrade} anfallarna. Dessa
\VUgras{drog sig tillbaka} och täckte \VUgras{slagfältet} med
\VUgras{döda} och \VUgras{sårade}. \VUgras{Sjukbärare} bar bort de sistnämnda
till \VUgras{ambulansen} för att låta \VUgras{kirurgen} vårda dem.

%%K t16-05-3 p
Trots det hjältemodiga motståndet från några \VUgras{bataljoner} som hade bildat en
\VUgras{karré}, var \VUgras{nederlaget} fullständigt, resten av armén
\VUgras{flydde}, och lämnade efter sig många \VUgras{fångar}. Fienden gick att
fullborda sin \VUgras{seger} genom att besätta staden.

%%K t16-05-3 p suite
Kanske kommer det att bli slutet på fälttåget och kommer det, efter ett
\VUgras{stillestånd}, att bli möjligt att underteckna ett \VUgras{fredsfördrag}.

%%K t16-tit-3 sub
\VUsaut{4.4mm}
\VUpk{10.2pt}{40}{Gåvor till nyåret.}
\VUsaut{2.8mm}

%%K t16-06-1 p
6. --- De unga soldaterna, som skrattade medan de lyssnade till veteranens pittoreska
berättelse bad honom om några förklaringar till. Vad mig beträffar, gick jag runt i boden för
att betrakta ett vackert \VUgras{armborst} \textsuperscript{(32)} och en \VUgras{båge}
\textsuperscript{(33)} med sin
\VUgras{pil} \textsuperscript{(31)}. Där fann jag en annan samling djur: två
\VUgras{sångfåglar} \textsuperscript{(34)}: en automatisk \VUgras{näktergal}
och en \VUgras{sångare} som sjöng av full hals, och en rad främmande djur: en
\VUgras{fladdermus} \textsuperscript{(35)}, en \VUgras{boaorm}
\textsuperscript{(36)}, en \VUgras{sköldpadda} \textsuperscript{(37)}, en
\VUgras{säl} \textsuperscript{(38)}, en \VUgras{val} \textsuperscript{(39)} och
en \VUgras{haj} \textsuperscript{(40)} av förgyllt läder.

%%K t16-07-1 p
7. --- Jag stannade inte länge för att betrakta dem, för jag började se det som jag drömde om
: en praktfull \VUgras{dockteater} \textsuperscript{(41)} med lysande \VUgras{dekorationer}
och en förtjusande röd \VUgras{ridå}. På
\VUgras{scenen} tycktes två skådespelare spela en \VUgras{roll} i ett mycket
fängslande \VUgras{skådespel}, för de små \VUgras{åskådarna} av papp som var installerade i
logerna eller satt på \VUgras{fåtöljerna} och på \VUgras{parketten} bakom \VUgras{dirigenten},
applåderade livligt.

%%K t16-07-2 p
Tyvärr kostade den teatern mycket.

%%K t16-08-1 p
8. --- För övrigt var det svårt att välja leksakerna, för i skyltfönstren var alla slags
leksaker upptravade: en \VUgras{Harlekin} \textsuperscript{(42)}, en
\VUgras{Pulcinella} \textsuperscript{(43)}, en \VUgras{Pierrot}
\textsuperscript{(44)}, \VUgras{ugglor} \textsuperscript{(45)} som slog med vingarna,
visslande \VUgras{koltrastar} \textsuperscript{(46)}, skällande
\VUgras{pudlar}, en \VUgras{fonograf} \textsuperscript{(47)} och
\VUgras{tålamodsspel}.

%%K t16-noto-1 noto
\VUnotes{143.2mm}{%
\textit{(*) Se tabell nr 6.}%
}

%%K t16-08-1 p suite
En av de leksakerna roade mig mycket: den föreställde ett \VUgras{sjöslag}
\textsuperscript{(48)} i vilket ett \VUgras{skepp} angreps av \VUgras{sjörövare} vid ett land
bevuxet med \VUgras{kokospalmer}.

%%K t16-09-1 p
9. --- Jag kunde mycket lätt betrakta alla de underverken, för det var bara få människor i
varuhuset, knappt några flanörer, en \VUgras{dragon} \textsuperscript{(49)} och en
\VUgras{inkasserare} \textsuperscript{(50)} som lämnade fram en växel vid huvud\VUgras{kassan}
\textsuperscript{(51)}.

%%K t16-10-1 p
10. --- Jag frågade min farfar vad de böcker används till som ligger på hyllor bakom
\VUgras{kassörskan}; han svarade mig att det är \VUgras{räkenskapsböcker}.

%%K t16-tit-4 sub
\VUsaut{3.6mm}
\VUpk{10.2pt}{40}{Att gapa vid en bod.}
\VUsaut{2.8mm}

%%K t16-11-1 p
11. --- När vi lämnade leksaksavdelningen, gick vi till sadelmakeriet för att kasta en blick
på \VUgras{sadlarna} \textsuperscript{(53)},
\VUgras{stigbyglarna} \textsuperscript{(54)}, \VUgras{tömmarna}
\textsuperscript{(55)}, \VUgras{betslen} \textsuperscript{(56)} och
\VUgras{ridpiskorna}. Medan \VUgras{avdelningschefen} \textsuperscript{(57)} gav
min farfar några upplysningar, gick jag och betraktade utställningen av
\VUgras{porslin} och \VUgras{kristall} \textsuperscript{(58)} där vackra
\VUgras{salladsskålar} och fina \VUgras{tekannor} \textsuperscript{(59)} står.

%%K t16-12-1 p
12. --- För att gå upp till första våningen, gick vi bakom skyltfönstret med
\VUgras{korsetter} \textsuperscript{(60)}, bredvid strumpavdelningen, där mycket
vackra \VUgras{underbyxor} \textsuperscript{(63)} var utställda. Alldeles intill lät en
\VUgras{försäljerska} \textsuperscript{(61)} en \VUgras{köperska} \textsuperscript{(62)} välja
vackra \VUgras{tyger} \textsuperscript{(64)} av siden.

%%K t16-12-2 p
Medan jag steg uppför trappan, märkte jag att varuhuset är prytt med japanska
\VUgras{lyktor} \textsuperscript{(65)}, \VUgras{parasoller}
\textsuperscript{(66)} och väldiga \VUgras{palmer} \textsuperscript{(70)}.

%%K t16-13-1 p
13. --- På första våningen fanns det inte mycket att se; bara möbler (*),
\VUgras{kassaskåp} \textsuperscript{(67)}, \VUgras{byråar}
\textsuperscript{(68)}, \VUgras{barnvagnar} \textsuperscript{(69)}. Men helhetssynen av
varuhuset var mycket \VUgras{förlustande}.

%%K t16-14-1 p
14. --- Vid våra fötter såg jag musikinstrumenten: \VUgras{harpor} \textsuperscript{(71)},
\VUgras{gitarrer} \textsuperscript{(72)},
\VUgras{mandoliner} \textsuperscript{(73)}, \VUgras{kornetter med ventiler}
\textsuperscript{(74)}, \VUgras{klarinetter} \textsuperscript{(75)}, violiner med sin
\VUgras{stråke} \textsuperscript{(76)} och till och med en \VUgras{stor trumma}
\textsuperscript{(77)}. Något längre bort, vid pappersavdelningen, prutade en \VUgras{student}
\textsuperscript{(78)} hos \VUgras{biträdet} \textsuperscript{(79)} på en skrivboksmapp.
Bredvid dem urskilde jag en vacker
\VUgras{abc-bok} \textsuperscript{(80)}.

%%K t16-14-2 p
Mitt i varuhuset var en hög \VUgras{julgran} \textsuperscript{(81)} rest som var helt försedd
med ljus, grannlåt och alla slags leksaker: \VUgras{trähästar} \textsuperscript{(82)},
\VUgras{dockor} \textsuperscript{(83)}, o.s.v..

%%K t16-14-3 p
\VUgras{Parfymhandeln} \textsuperscript{(84)} med sina \VUgras{tandmedel} och
olika \VUgras{parfymer} spred en mycket god doft som kom emot oss.

%%K t16-tit-5 sub
\VUsaut{3.4mm}
\VUpk{10.2pt}{40}{Vi köper något.}
\VUsaut{2.8mm}

%%K t16-15-1 p
15. --- Eftersom min farfar började tycka att tiden blev lång, gick vi ner för den stora
trappan. Han stannade ett ögonblick framför en \VUgras{astronomisk tabell}
\textsuperscript{(85)} som, sade han mig, föreställer \VUgras{kometer} \textsuperscript{(a)},
\VUgras{mån- och solförmörkelser} \textsuperscript{(b)}, en vindros med de
\VUgras{fyra väderstrecken} \textsuperscript{(c)}
\VUgras{(Norr, Söder, Öster och Väster)}, en karta över de två
\VUgras{halvkloten} \textsuperscript{(d)}, \VUgras{planeterna}
\textsuperscript{(e)} och \VUgras{solsystemet} \textsuperscript{(f)}. Jag förstod ingenting av
allt det, och jag blev mycket mer fängslad av den trensprydda livrén hos en
\VUgras{springpojke} \textsuperscript{(86)} som gick förbi, och av de vackra
\VUgras{solfjädrarna} \textsuperscript{(87)} som låg bredvid mig, vid \VUgras{portmonnäerna}
\textsuperscript{(88)} och
\VUgras{plånböckerna} \textsuperscript{(89)}.

%%K t16-16-1 p
16. --- Jag beundrade också på skyltningen \VUgras{fjädrar} \textsuperscript{(90)}, och
\VUgras{bältesspännen} \textsuperscript{(91)}, de vackra \VUgras{konstgjorda blommorna}
\textsuperscript{(94)} som var graciöst ordnade i korgar.

%%K t16-16-1 p suite
\VUgras{Violerna}, \VUgras{lövkojorna}, \VUgras{hyacinterna}
\textsuperscript{(95)}, \VUgras{pelargonierna} \textsuperscript{(96)},
\VUgras{liljorna} \textsuperscript{(97)} och \VUgras{krasserna}
\textsuperscript{(98)} var mycket väl efterhärmade.

%%K t16-tit-6 sub
\VUsaut{4.4mm}
\VUpk{10.2pt}{40}{En gåva av min farfar.}
\VUsaut{2.8mm}

%%K t16-17-1 p
17. --- Jag bad min farfar att köpa åt mig en av de \VUgras{tandborstar}
\textsuperscript{(92)} som låg i en ask bredvid \VUgras{hårnålar} \textsuperscript{(93)}. Han
samtyckte och jag gick själv och betalade mitt inköp vid kassan, bredvid vilken man gjorde i
ordning en viktig \VUgras{sändning} \textsuperscript{(100)} åt en \VUgras{kund}
\textsuperscript{(99)}.

%%K t16-18-1 p
18. --- Plötsligt uppstod det en lätt uppståndelse bland de människor som hade stannat vid de
konstnärliga \VUgras{bronsarbetena} \textsuperscript{(101)}. En
\VUgras{tjuv} \textsuperscript{(102)} hade just tagits på bar gärning av en
\VUgras{inspektör} \textsuperscript{(103)}, medan han försökte bestjäla någon.
Man sörjde för att låta hämta en poliskonstapel för att \VUgras{häkta} honom, och just när vi
gick ut, fördes gärningsmannen till fängelset.

\VUsaut{6.0mm}
\VUfilet{22mm}
